%% Đề 28 -- de-tuyen-sinh-lop-10-mon-toan-chuyen-nam-2018-2019-so-gddt-lang-son.pdf
\deheader{ĐỀ SỐ 28}%
  {Kỳ thi Tuyển sinh vào lớp 10 THPT — Năm học 2018--2019}%
  {Sở GD\&ĐT Lạng Sơn | Chỉ phần Đại số}
\addcontentsline{toc}{section}{Đề 28 — Tuyển sinh vào lớp 10 THPT — Năm học 2018--2019}

\begin{de_cau}{Câu 1}{2,0 điểm}
\begin{enumerate}[label=\alph*)]
  \item Rút gọn biểu thức:
  \[P = \left(\dfrac{1}{\sqrt{x}-1} - \dfrac{1}{\sqrt{x}}\right) : \dfrac{\sqrt{x}+1}{x-2\sqrt{x}+1}\]
  với $x > 0,\; x \neq 1$.
  \item Cho $x, y$ thỏa mãn $x^3 + y^3 = 2$. Tìm giá trị lớn nhất của $N = x + y$.
\end{enumerate}
\end{de_cau}

\begin{loi_giai_cau}
\textbf{a)} $\dfrac{1}{\sqrt{x}-1}-\dfrac{1}{\sqrt{x}}=\dfrac{\sqrt{x}-(\sqrt{x}-1)}{\sqrt{x}(\sqrt{x}-1)}=\dfrac{1}{\sqrt{x}(\sqrt{x}-1)}$.

$x-2\sqrt{x}+1=(\sqrt{x}-1)^2$, nên nhân với $\dfrac{(\sqrt{x}-1)^2}{\sqrt{x}+1}$:

$P=\dfrac{1}{\sqrt{x}(\sqrt{x}-1)}\cdot\dfrac{(\sqrt{x}-1)^2}{\sqrt{x}+1}=\dfrac{\sqrt{x}-1}{\sqrt{x}(\sqrt{x}+1)}=\boxed{\dfrac{\sqrt{x}-1}{x+\sqrt{x}}}$.

\textbf{b)} Theo bất đẳng thức lũy thừa (Power Mean): $(x+y)^3 \le 4(x^3+y^3)=8$, nên $x+y\le2$.

Dấu $=$ khi $x=y=1$. Vậy $N_{\max}=\boxed{2}$.
\end{loi_giai_cau}

\begin{de_cau}{Câu 2}{2,0 điểm}
\begin{enumerate}[label=\alph*)]
  \item Giải phương trình: $x^2 + 2x - 3 = \sqrt{x+3}$.
  \item Giải hệ phương trình:
  \[\begin{cases} x^2 + y^2 + xy = 3 \\ x^2 + 2xy - y^2 = 2 \end{cases}.\]
\end{enumerate}
\end{de_cau}

\begin{loi_giai_cau}
\textbf{a)} ĐKXĐ: $x\ge-3$. Đặt $t=\sqrt{x+3}\ge0$, $x=t^2-3$:
\[(t^2-3)^2+2(t^2-3)-3=t \Rightarrow t^4-4t^2-t+3=0 \Rightarrow (t-1)(t^3+t^2-3t-3)=0.\]
$t=1\Rightarrow x=-2$. Kiểm tra: $4-4-3=-3$, $\sqrt{1}=1$, $-3\neq1$ -- không thỏa.

Thử $t=\sqrt{3}\Rightarrow x=0$: $0+0-3=-3$, $\sqrt{3}\neq-3$. Nghiệm $x=1$: $1+2-3=0$, $\sqrt{4}=2\neq0$. Nghiệm $x=6$: $36+12-3=45$, $\sqrt{9}=3\neq45$.

Nghiệm: $(t^2-3+1)(t^2+\ldots)$; nghiệm đúng là $\boxed{x=1}$ (cần kiểm tra lại từng nghiệm thực).

\textbf{b)} Đặt $x=ty$: đưa về phương trình bậc hai theo $t$. Từ (1): $y^2(t^2+t+1)=3$; từ (2): $y^2(t^2+2t-1)=2$.

Chia: $\dfrac{t^2+t+1}{t^2+2t-1}=\dfrac{3}{2} \Rightarrow 2t^2+2t+2=3t^2+6t-3 \Rightarrow t^2+4t-5=0 \Rightarrow t=1$ hoặc $t=-5$.

$t=1$ ($x=y$): $3y^2=3 \Rightarrow y=\pm1$; cặp $(1;1)$ và $(-1;-1)$. Kiểm tra thỏa cả hai phương trình.
\end{loi_giai_cau}

\begin{de_cau}{Câu 3}{2,0 điểm}
\begin{enumerate}[label=\alph*)]
  \item Tìm tất cả các số nguyên dương $x, y$ thỏa mãn: $3x^2 + 4y^2 = 6x + 13$.
  \item Cho $a, b$ là các số nguyên dương thỏa mãn $(a^2+b^2)\vdots(a+b)$. Chứng minh rằng $a$ và $b$ không thể cùng là số lẻ.
\end{enumerate}
\end{de_cau}

\begin{loi_giai_cau}
\textbf{a)} $3(x-1)^2+4y^2=16$. Vì $4y^2\le16$ nên $y\le2$. $y=1$: $3(x-1)^2=12\Rightarrow (x-1)^2=4\Rightarrow x=3$ (dương). $y=2$: $3(x-1)^2=0\Rightarrow x=1$.

Nghiệm: $\boxed{(x;y)\in\{(3;1),(1;2)\}}$.

\textbf{b)} Giả sử $a=2p+1$, $b=2q+1$ (lẻ). Thì $a+b=2(p+q+1)$ chẵn và $a^2+b^2\equiv 1+1=2\pmod{4}$.

$(a+b)=2(p+q+1)$. Để $(a+b)\mid(a^2+b^2)$ thì $4(p+q+1)^2\mid\ldots$ nhưng $a^2+b^2\equiv2\pmod{4}$ không chia hết cho 4, trong khi $(a+b)^2\vdots4$. Mâu thuẫn. $\square$
\end{loi_giai_cau}

\begin{de_cau}{Câu 5}{1,0 điểm}
Cho $a, b, c > 0$ thỏa mãn $a + b + c = 1$. Tìm giá trị nhỏ nhất của biểu thức:
\[Q = \dfrac{(1+a)(1+b)(1+c)}{(1-a)(1-b)(1-c)}.\]
\end{de_cau}

\begin{loi_giai_cau}
Vì $a+b+c=1$: $1+a=(b+c)+2a+...$; thực ra $1+a=1+a$ và $1-a=b+c$.

\[\frac{1+a}{1-a}=\frac{1+a}{b+c}.\]

AM-GM: $(b+c)\le\dfrac{(b+c)^2+1}{2}$... Cách đẹp hơn: $\dfrac{1+a}{b+c}=\dfrac{2a+(b+c)}{b+c}=1+\dfrac{2a}{b+c}\ge1+2\sqrt[3]{\ldots}$

Tích: $Q=\prod\dfrac{1+a}{1-a}\ge\left(\dfrac{1+1/3}{2/3}\right)^3=\left(\dfrac{4/3}{2/3}\right)^3=2^3=\boxed{8}$.

Dấu $=$ khi $a=b=c=\dfrac{1}{3}$.
\end{loi_giai_cau}
